% ══════════════════════════════════════════════════════════════════
% Abstract
% ══════════════════════════════════════════════════════════════════

Neural network surrogates for hyperelastic strain energy functions offer a
promising path toward accelerating nonlinear finite element simulations, yet
practical deployment remains hindered by (i) the lack of systematic comparisons
across network architectures and material classes, (ii) the gap between
Python-trained models and production Fortran solvers, and (iii) insufficient
validation against reference finite element solutions.  This work addresses all
three challenges through \texttt{hyper-surrogate}, an open-source framework
that provides a complete pipeline from data generation through training to
Fortran~90 code emission.

We benchmark three neural architectures --- multilayer perceptron (MLP),
input-convex neural network (ICNN), and polyconvex ICNN --- on five
hyperelastic materials spanning isotropic (Neo-Hookean, Mooney--Rivlin, Yeoh)
and anisotropic fiber-reinforced (Holzapfel--Ogden, Gasser--Ogden--Holzapfel)
formulations.  All models are trained with a physics-informed loss that
simultaneously matches strain energy and its invariant gradients via automatic
differentiation.  Accuracy, convergence, data efficiency, and inference speed
are compared quantitatively.

The trained networks are exported as self-contained Fortran~90 UMAT
subroutines, where the neural network evaluates the strain energy while
analytical kinematics compute stress and the consistent tangent operator.
Single-element finite element tests in Abaqus and FEAP demonstrate that the
hybrid NN--UMAT reproduces analytical Cauchy stress to within
% TODO: Insert actual tolerance from FE results.
\textbf{[XX]\%} relative error under uniaxial, biaxial, and shear loading.

% TODO: Add key quantitative results (best R^2, speedup factor, etc.)
% once benchmarks are finalized.
